\scanpage{15}
\textbf{1600 — Série A}\hfill\textbf{C. R. Acad. Sc. Paris, t. 278 (24 juin 1974)}

\rule{\linewidth}{0.4pt}

\textsc{Proposition 1}. — \emph{L'application $\alpha$ est un isomorphisme et nous avons la suite exacte} :

\diagram{\begin{tikzcd}[ampersand replacement=\&]
0\arrow[r]\&(8\mathbf Z)^r\arrow[r]\&\operatorname{Witt}(\mathrm K)\arrow[r]\&\operatorname{HW}(\mathrm K)\arrow[r]\&0
\end{tikzcd}}

\emph{où $r$ est le nombre de valuations réelles de} K.

\emph{Preuve}. — L'application $\beta$ est injective d'après le théorème de Hasse-Minkowski et
il en est de même de l'application $\alpha$. D'autre part, si $v$ est une valuation du type (i),
$\operatorname{LW}(\mathrm K_v)\cong8\mathbf Z$, et si elle est du type (ii), $\operatorname{LW}(\mathrm K_v)=\{0\}$. Par conséquent,
$\displaystyle\prod_{v\in\Omega}\operatorname{LW}(\mathrm K)\cong(8\mathbf Z)^r$. Le fait que $\alpha$ soit surjective résulte de la théorie classique des formes
quadratiques sur un corps de nombres [\emph{cf.} (${}^3$), p. 413 par exemple]\quad$\blacksquare$

D'après (${}^6$) ou (${}^7$), nous avons la suite exacte :

\diagram{\begin{tikzcd}[ampersand replacement=\&]
0\arrow[r]\&\mathrm N(\mathrm K)\arrow[r]\&\operatorname{HW}_0(\mathrm K)\arrow[r]\&\mathrm K^*/\mathrm K^{*2}\arrow[r]\&0
\end{tikzcd}}

où $\mathrm N(\mathrm K)$ est un sous-groupe de $\operatorname{Br}(\mathrm K)$ constitué d'éléments d'exponent deux de $\operatorname{Br}(\mathrm K)$
et $\operatorname{HW}_0(\mathrm K)$ l'image par $f_{\mathrm K}$ du sous-groupe $\operatorname{Witt}_0(\mathrm K)$ obtenu à partir des formes quadratiques
de rang pair.

On vérifie sans peine à l'aide de la s. e. [\emph{cf.} (${}^{10}$)] :

\diagram{\begin{tikzcd}[ampersand replacement=\&]
0\arrow[r]\&\operatorname{Br}(\mathrm K)\arrow[r]\&
\displaystyle\unsure{\bigoplus|\bigotimes}_{v\in\Omega}\operatorname{Br}(\mathrm K_v)\arrow[r]\&\mathbf Q/\mathbf Z\arrow[r]\&0
\end{tikzcd}}

et du résultat rappelé dans la démonstration de la proposition 1 que $\mathrm N(\mathrm K)=\operatorname{Br}^2(\mathrm K)$.
Le calcul de $\operatorname{Witt}(\mathrm K)$ est bien ramené à celui de $\mathrm K^*/\mathrm K^{*2}$.

3. \textsc{Les anneaux d'entiers algébriques}. — Soient K un corps de nombres et A l'anneau
des éléments de K entiers sur Z dont le corps de fractions est K.

Nous allons étudier le diagramme

\diagram{\begin{tikzcd}[ampersand replacement=\&,column sep=small]
0\arrow[r]\&\operatorname{LW}(\mathrm A)\arrow[r]\arrow[d,"a"']\&\operatorname{Witt}(\mathrm A)\arrow[r]\arrow[d,"b"']\&\operatorname{HW}(\mathrm A)\arrow[r]\arrow[d,"c"']\&0\\
0\arrow[r]\&\operatorname{LW}(\mathrm K)\arrow[r]\&\operatorname{Witt}(\mathrm K)\arrow[r]\&\operatorname{HW}(\mathrm K)\arrow[r]\&0
\end{tikzcd}}

\textsc{Lemme 1}. — \emph{L'application $c$ de $\operatorname{HW}(\mathrm A)$ dans $\operatorname{HW}(\mathrm K)$ est injective.}

\emph{Preuve}. — Considérons les diagrammes

\diagram{\begin{tikzcd}[ampersand replacement=\&,column sep=small]
0\arrow[r]\&\mathrm N(\mathrm A)\arrow[r]\arrow[d]\&\operatorname{HW}_0(\mathrm A)\arrow[r]\arrow[d]\&\mathrm Q(\mathrm A)\arrow[r]\arrow[d]\&0\\
0\arrow[r]\&\mathrm N(\mathrm K)\arrow[r]\&\operatorname{HW}_0(\mathrm K)\arrow[r]\&\mathrm Q(\mathrm K)\arrow[r]\&0
\end{tikzcd}}

et

\diagram{\begin{tikzcd}[ampersand replacement=\&]
0\arrow[r]\&\mathrm N(\mathrm A)\arrow[r]\arrow[d]\&\operatorname{Br}(\mathrm A)\arrow[d]\\
0\arrow[r]\&\mathrm N(\mathrm K)\arrow[r]\&\operatorname{Br}(\mathrm K)
\end{tikzcd}}

Il est montré dans (${}^1$) que $\operatorname{Br}(\mathrm A)$ est un sous-groupe de $\operatorname{Br}(\mathrm K)$ car A est un anneau de
Dedekind. D'autre part, le groupe $\mathrm Q(\mathrm R)$ d'un anneau quelconque R est identique au groupe

\mspage{14}
